%! The Rational Parent Function \& Transformations — Unit 4, Lesson 4.4 — Student Packet Cover Sheet
\documentclass[10pt]{article}
\usepackage{algebra2-article}
\usepackage{algebra2-boxes}
\usepackage{ltablex}
\keepXColumns

\begin{document}

% Full-bleed forest banner
\begin{tikzpicture}[remember picture, overlay]
  \fill[forest] (current page.north west)
    rectangle ([yshift=-0.9in]current page.north east);
\end{tikzpicture}
\vspace{-0.8in}

\begin{center}
  {\color{white}\textbf{\LARGE Algebra 2}}\\[4pt]
  {\color{forestmid}\large\textbf{Unit 4 \quad Rational Functions}}\\[6pt]
  {\color{white}\textbf{\large Lesson 4.4 \quad The Rational Parent Function \& Transformations}}
\end{center}

\vspace{0.22in}
\namedateperiod
\vspace{0.18in}

\begin{learningtargetbox}
\small
\begin{enumerate}[leftmargin=1.5em, itemsep=5pt]
  \item \ldots{} graph and describe the \textbf{rational parent function} $f(x)=\frac1x$ --- a
        \textbf{hyperbola} with two branches, a vertical asymptote $x=0$, a horizontal asymptote
        $y=0$, origin symmetry, and \emph{no intercepts at all}.
  \item \ldots{} read $a$, $h$, and $k$ out of $g(x)=\dfrac{a}{x-h}+k$ and use them to place the
        \textbf{asymptotes first} --- $x=h$ from the denominator (which does the opposite of what it
        says) and $y=k$ from the number outside --- then the domain, the range, and the intercepts.
  \item \ldots{} \textbf{write the equation} of a rational function from its graph: the asymptotes
        give me $h$ and $k$, and one point on the curve gives me $a$.
  \item \ldots{} compare the rational parent with the parents I already know, and recognize a
        function whose top and bottom are both \textbf{linear} as a disguised $\frac1x$ by rewriting
        its numerator with its denominator.
\end{enumerate}
\end{learningtargetbox}

\vspace{0.14in}

\begin{tocbox}
\small
\renewcommand{\arraystretch}{1.5}
\begin{tabularx}{\linewidth}{@{} c l X r @{}}
\rowcolor{forestbg}
\textbf{\#} & \textbf{Component} & \textbf{Description} & \textbf{Score} \\
\midrule
1 & Warm-Up        & Spiral review: the reciprocal table, the shift rules from Units 1--2, and one rational function written two ways & \blank{1.2cm} \\
2 & Guided Notes   & The parent $y=\frac1x$, the two shifts that move the asymptotes, the stretch, the general form, and writing an equation from a graph & \blank{1.2cm} \\
3 & Group Activity & \emph{Move the Walls} --- asymptote drills, a graph-matching set, an error analysis, the disguise, and a dilution model, by tier & \blank{1.2cm} \\
4 & Exit Ticket    & Quick check: read one equation completely, write one equation from a graph, plus an SOL-style item & \blank{1.2cm} \\
5 & Homework       & Asymptotes at a glance, full feature reads, equations from graphs, four disguises, a table comparison, and a free-throw model & \blank{1.2cm} \\
\midrule
  & \multicolumn{2}{r}{\textbf{Total}} & \blank{1.2cm} \\
\end{tabularx}
\end{tocbox}

\vspace{0.10in}

\begin{remindbox}
\small The \textbf{rational parent function} is $f(x)=\frac1x$. Its graph is a \textbf{hyperbola}: two
separate \textbf{branches}, a vertical asymptote at $x=0$, a horizontal asymptote at $y=0$, and --- unlike
every parent before it --- \textbf{no $x$-intercept and no $y$-intercept}, because a fraction is zero only
when its numerator is and $x=0$ is not in the domain. Every transformation of it can be written
$g(x)=\dfrac{a}{x-h}+k$, and the one sentence to carry into every problem is this: \textbf{the asymptotes
are the parent's axes, and they travel with the graph}. Set the denominator to zero for the vertical
asymptote $x=h$; the number added outside is the horizontal asymptote $y=k$; the two cross at the
\textbf{center} $(h,k)$, which is \emph{not} a point of the graph. The number inside the denominator does
the \emph{opposite} of what it says --- $x-3$ shifts \emph{right} $3$ --- exactly as it did for
$|x-3|$ and $(x-3)^2$. The value of $a$ moves no asymptote at all: it pushes the branches away from the
center when $|a|>1$, pulls them in when $|a|<1$, and flips them to the opposite pair of regions when
$a<0$. The domain excludes $h$ and the range excludes $k$; both intercepts are \emph{computed}, by
setting $x=0$ and by solving $g(x)=0$. To go the other way --- from a graph to an equation --- read $h$
and $k$ off the asymptotes, substitute one clearly marked point, and solve for $a$; then check your
equation on a second point. Finally, a function whose numerator and denominator are both \textbf{linear}
is one of these in disguise: rewrite the numerator in terms of the denominator (as in
$\frac{x-1}{x-3}=\frac{(x-3)+2}{x-3}=1+\frac{2}{x-3}$) and the asymptotes appear.
\end{remindbox}

\end{document}
